\documentclass[a4paper,11pt]{article}
\usepackage[utf8]{inputenc}
\usepackage[T1]{fontenc}
\usepackage[french]{babel}
\usepackage{amsmath,amssymb}
\usepackage{geometry}
\usepackage{tikz}
\usetikzlibrary{positioning,shapes.geometric,calc,arrows.meta}
\usepackage{listings}
\usepackage{booktabs}
\usepackage{xcolor}
\usepackage{array}
\usepackage{mdframed}
\usepackage{enumitem}
\usepackage{tabularx}

\geometry{top=2cm, bottom=2cm, left=2.5cm, right=2.5cm}

% Couleurs
\definecolor{maxblue}{RGB}{41,128,185}
\definecolor{minred}{RGB}{192,57,43}
\definecolor{gris}{RGB}{160,160,160}
\definecolor{vertok}{RGB}{39,174,96}
\definecolor{bgcode}{RGB}{245,245,245}

% Styles tikz
\tikzset{
  abr/.style={circle, draw=black, thick, minimum size=8mm, font=\small\bfseries, fill=white},
  abrh/.style={circle, draw=black, thick, minimum size=8mm, font=\tiny, fill=white},
  maxnode/.style={circle, draw=maxblue, fill=maxblue!15, thick,
                  minimum size=9mm, font=\small\bfseries},
  minnode/.style={circle, draw=minred, fill=minred!15, thick,
                  minimum size=9mm, font=\small\bfseries},
  leafnode/.style={rectangle, draw=black, fill=gray!15, thick,
                   minimum size=8mm, font=\small\bfseries},
  prunedleaf/.style={rectangle, draw=gris, fill=gray!5, thick,
                     minimum size=8mm, font=\small, text=gris, dashed},
  prumednode/.style={circle, draw=gris, fill=gray!5, thick,
                     minimum size=9mm, font=\small, text=gris, dashed},
  edge from parent/.style={draw, thick, -},
  level distance=1.4cm,
}

% Style listings
\lstset{
  basicstyle=\small\ttfamily,
  backgroundcolor=\color{bgcode},
  frame=single,
  framerule=0.4pt,
  breaklines=true,
  language=C,
  keywordstyle=\color{blue}\bfseries,
  commentstyle=\color{gray}\itshape,
  numbers=none,
}

\newcommand{\boitereponse}[1]{%
  \begin{mdframed}[linecolor=black!40, backgroundcolor=white, linewidth=0.5pt,
                   innertopmargin=6pt, innerbottommargin=6pt]
  #1
  \end{mdframed}
}

\setlist[itemize]{topsep=4pt, itemsep=2pt}
\setlist[enumerate]{topsep=4pt, itemsep=2pt}

% =====================================================================
\begin{document}

\begin{center}
  {\LARGE\bfseries Correction --- TD Complexité Algorithmique}\\[0.4em]
  {\large Exercices 2 \& 3 \textemdash{} ABR/AVL et Minimax/Alpha-bêta}\\[0.2em]
  {\small IMT Mines Alès --- Module 8.2.3}
\end{center}
\rule{\linewidth}{0.8pt}

% =====================================================================
\section*{Exercice 2 --- Structures équilibrées et coût mémoire}
% =====================================================================

\subsection*{Partie A --- ABR et coût de la recherche}

%------------------------------------------------------------
\paragraph{Question 1 --- ABR pour 42, 17, 63, 8, 31, 55, 78.}

On insère les clés dans l'ordre : la première devient la racine, chaque nouvelle clé descend à gauche si elle est plus petite, à droite sinon.

\begin{center}
\begin{tikzpicture}[
  level 1/.style={sibling distance=4.5cm},
  level 2/.style={sibling distance=2.2cm},
]
  \node[abr]{42}
    child{node[abr]{17}
      child{node[abr]{8}}
      child{node[abr]{31}}
    }
    child{node[abr]{63}
      child{node[abr]{55}}
      child{node[abr]{78}}
    };
\end{tikzpicture}
\end{center}

\boitereponse{%
\textbf{Hauteur} (nombre d'arêtes de la racine jusqu'à la feuille la plus profonde) $= \mathbf{2}$.%
}

%------------------------------------------------------------
\paragraph{Question 2 --- Coût de la recherche.}

Avec 7 nœuds répartis sur 3 niveaux (1 nœud niveau 1, 2 nœuds niveau 2, 4 feuilles niveau 3) :

\begin{itemize}
  \item \textbf{Pire cas} : la clé cherchée est une feuille $\to$ \textbf{3 comparaisons} ($= h + 1$).
  \item \textbf{Cas moyen} (en supposant chaque nœud équiprobable) :
  \[
    \bar{c} = \frac{1 \times 1 + 2 \times 2 + 4 \times 3}{7}
             = \frac{1 + 4 + 12}{7}
             = \frac{17}{7}
             \approx \mathbf{2{,}43} \textbf{ comparaisons.}
  \]
\end{itemize}

%------------------------------------------------------------
\paragraph{Question 3 --- ABR dégénéré (insertion triée).}

En insérant dans l'ordre croissant 8, 17, 31, 42, 55, 63, 78, chaque clé est systématiquement plus grande que la précédente : l'arbre dégénère en liste chaînée vers la droite.

\begin{center}
\begin{tikzpicture}[node distance=0pt]
  \node[abr] (n1) {8};
  \node[abr, right=2.5cm of n1, yshift=-1.2cm] (n2) {17};
  \node[abr, right=2.5cm of n2, yshift=-1.2cm] (n3) {31};
  \node[abr, right=2.5cm of n3, yshift=-1.2cm] (n4) {42};
  \node[abr, right=2.5cm of n4, yshift=-1.2cm] (n5) {55};
  \node[abr, right=2.5cm of n5, yshift=-1.2cm] (n6) {63};
  \node[abr, right=2.5cm of n6, yshift=-1.2cm] (n7) {78};
  % petits tirets pour indiquer pas de fils gauche
  \draw[thick] (n1) -- (n2) node[midway,above,sloped,font=\tiny]{droite};
  \draw[thick] (n2) -- (n3) node[midway,above,sloped,font=\tiny]{droite};
  \draw[thick] (n3) -- (n4) node[midway,above,sloped,font=\tiny]{droite};
  \draw[thick] (n4) -- (n5) node[midway,above,sloped,font=\tiny]{droite};
  \draw[thick] (n5) -- (n6) node[midway,above,sloped,font=\tiny]{droite};
  \draw[thick] (n6) -- (n7) node[midway,above,sloped,font=\tiny]{droite};
\end{tikzpicture}
\end{center}

\boitereponse{%
\textbf{Hauteur} $= n - 1 = \mathbf{6}$.\\
\textbf{Pire cas} : chercher 78 $\to$ \textbf{7 comparaisons} ($= n$).%
}

%------------------------------------------------------------
\paragraph{Question 4 --- Conclusion.}

\boitereponse{%
L'ordre d'insertion détermine directement la forme de l'ABR.
Un ordre trié (croissant ou décroissant) produit un arbre \emph{dégénéré} de hauteur $O(n)$ et un coût de recherche $O(n)$.
Un ordre aléatoire donne en moyenne une hauteur $O(\log n)$.
\textbf{Sans mécanisme d'équilibrage, les performances d'un ABR sont donc imprévisibles.}%
}

% ============================
\subsection*{Partie B --- Rotations et équilibrage AVL}

On rappelle la convention : le \textbf{facteur d'équilibre} (bf) d'un nœud $= h(\text{sous-arbre gauche}) - h(\text{sous-arbre droit})$, avec $h(\text{null}) = -1$.

%------------------------------------------------------------
\paragraph{Question 1 --- Insertion de la clé 5.}

$5 < 42 \to \text{gauche}$, $5 < 17 \to \text{gauche}$, $5 < 8 \to$ fils gauche de 8.

\begin{center}
\begin{tikzpicture}[
  level 1/.style={sibling distance=4.5cm},
  level 2/.style={sibling distance=2.2cm},
  level 3/.style={sibling distance=1.5cm},
]
  \node[abr]{42}
    child{node[abr]{17}
      child{node[abr]{8}
        child{node[abr]{5}}
        child[missing]{}
      }
      child{node[abr]{31}}
    }
    child{node[abr]{63}
      child{node[abr]{55}}
      child{node[abr]{78}}
    };
\end{tikzpicture}
\end{center}

Facteurs d'équilibre en remontant de la feuille vers la racine :

\begin{center}
\begin{tabular}{cccc}
\toprule
Nœud & $h_g$ & $h_d$ & bf \\
\midrule
5    & $-1$ & $-1$ & 0 \\
8    & 0    & $-1$ & \textbf{1} \\
31   & $-1$ & $-1$ & 0 \\
17   & 1    & 0    & \textbf{1} \\
55, 78 & \multicolumn{3}{c}{0} \\
63   & 0    & 0    & 0 \\
42   & 2    & 1    & \textbf{1} \\
\bottomrule
\end{tabular}
\end{center}

\boitereponse{Tous les facteurs sont dans $\{-1, 0, 1\}$ : \textbf{l'arbre est toujours AVL}.}

%------------------------------------------------------------
\paragraph{Question 2 --- Insertion de la clé 3.}

$3 < 42 \to 17 \to 8 \to 5 \to$ fils gauche de 5.

\begin{center}
\begin{tikzpicture}[
  level 1/.style={sibling distance=4.5cm},
  level 2/.style={sibling distance=2.2cm},
  level 3/.style={sibling distance=1.5cm},
  level 4/.style={sibling distance=1cm},
]
  \node[abr]{42}
    child{node[abr]{17}
      child{node[abr, label={[font=\small,text=red]right:bf=\textbf{2}!}]{8}
        child{node[abr]{5}
          child{node[abr]{3}}
          child[missing]{}
        }
        child[missing]{}
      }
      child{node[abr]{31}}
    }
    child{node[abr]{63}
      child{node[abr]{55}}
      child{node[abr]{78}}
    };
\end{tikzpicture}
\end{center}

\begin{itemize}
  \item $h_g(8) = h(5) = 1$, $h_d(8) = -1$ $\Rightarrow$ bf(8) $= 2$ : \textbf{nœud déséquilibré = 8}.
  \item Configuration : 3 est dans le sous-arbre \emph{gauche du fils gauche} de 8 $\to$ cas \textbf{Gauche-Gauche}.
  \item Rotation nécessaire : \textbf{rotation simple droite} sur le nœud 8.
\end{itemize}

Application de \texttt{rotate\_right(8)} : $y = 8$, $x = 5$, $T_2 = \texttt{null}$ (fils droit de 5).

\begin{center}
\begin{tikzpicture}[
  level 1/.style={sibling distance=4.5cm},
  level 2/.style={sibling distance=2.2cm},
  level 3/.style={sibling distance=1.5cm},
]
  \node[abr]{42}
    child{node[abr]{17}
      child{node[abr]{5}
        child{node[abr]{3}}
        child{node[abr]{8}}
      }
      child{node[abr]{31}}
    }
    child{node[abr]{63}
      child{node[abr]{55}}
      child{node[abr]{78}}
    };
\end{tikzpicture}
\end{center}

\boitereponse{Facteurs d'équilibre après rotation : 3:0, 8:0, 5:0, 31:0, 17:1, 55:0, 78:0, 63:0, 42:1. \\
Tous dans $\{-1,0,1\}$ : \textbf{arbre AVL valide}.}

%------------------------------------------------------------
\paragraph{Question 3 --- Construction AVL avec 8, 17, 31, 42, 55, 63, 78.}

On trace les étapes clés (toutes les rotations sont des \emph{rotations simples gauches} car les insertions suivent toujours le sous-arbre droit) :

\begin{enumerate}[label=\textbf{Étape \arabic*.}]
  \item Insère \textbf{8, 17} : bf(8) $= -1$ \checkmark
  \item Insère \textbf{31} : bf(8) $= -2$ (RR) $\to$ \textbf{rotation gauche sur 8}
        $\Rightarrow$ racine = 17, fils : \{8, 31\}
  \item Insère \textbf{42} : bf(31) $= -1$, bf(17) $= -1$ \checkmark
  \item Insère \textbf{55} : bf(31) $= -2$ (RR) $\to$ \textbf{rotation gauche sur 31}
        $\Rightarrow$ sous-arbre droit de 17 devient \{42, g:31, d:55\}
  \item Insère \textbf{63} : bf(17) $= -2$ (RR) $\to$ \textbf{rotation gauche sur 17 (la racine)}
        $\Rightarrow$ racine = 42, g:\{17,g:8,d:31\}, d:\{55,d:63\}
  \item Insère \textbf{78} : bf(55) $= -2$ (RR) $\to$ \textbf{rotation gauche sur 55}
        $\Rightarrow$ sous-arbre droit de 42 devient \{63, g:55, d:78\}
\end{enumerate}

\textbf{Arbre AVL final :}

\begin{center}
\begin{tikzpicture}[
  level 1/.style={sibling distance=4.5cm},
  level 2/.style={sibling distance=2.2cm},
]
  \node[abr]{42}
    child{node[abr]{17}
      child{node[abr]{8}}
      child{node[abr]{31}}
    }
    child{node[abr]{63}
      child{node[abr]{55}}
      child{node[abr]{78}}
    };
\end{tikzpicture}
\end{center}

\boitereponse{%
\textbf{Hauteur AVL} $= 2$.\\
\textbf{Hauteur ABR dégénéré} $= 6$.\\
L'AVL a reconstruit un arbre parfaitement équilibré, identique à celui de la Question~1, quelle que soit l'ordre d'insertion.%
}

%------------------------------------------------------------
\paragraph{Question 4 --- Coût $O(\log n)$ d'une insertion AVL.}

\boitereponse{%
\begin{itemize}
  \item \textbf{Descente} : on parcourt au plus $h = O(\log n)$ nœuds pour trouver la position d'insertion.
  \item \textbf{Rééquilibrage} : on remonte jusqu'à la racine en recalculant les facteurs d'équilibre, soit $O(\log n)$ nœuds. Chaque déséquilibre se corrige avec \textbf{au plus une rotation simple ou double}, coût $O(1)$.
  \item \textbf{Coût total} : $O(\log n) + O(1) = \mathbf{O(\log n)}$.
\end{itemize}%
}

% ============================
\subsection*{Partie C --- Application : table de transposition}

%------------------------------------------------------------
\paragraph{Question 1 --- Taille d'un nœud \texttt{AvlNode}.}

\begin{center}
\begin{tabular}{lcc}
\toprule
Champ         & Type            & Taille \\
\midrule
\texttt{key}   & \texttt{uint64\_t} & 8 octets \\
\texttt{value} & \texttt{int}       & 4 octets \\
\texttt{height}& \texttt{int}       & 4 octets \\
\texttt{*left} & pointeur 64 bits   & 8 octets \\
\texttt{*right}& pointeur 64 bits   & 8 octets \\
\midrule
\textbf{Total} &                    & \textbf{32 octets} \\
\bottomrule
\end{tabular}
\end{center}

Pas de padding : \texttt{key}(8) + \texttt{value}(4) + \texttt{height}(4) + \texttt{left}(8) + \texttt{right}(8) = 32 octets, tous les champs sont naturellement alignés.

%------------------------------------------------------------
\paragraph{Question 2 --- Mémoire pour $10^6$ positions.}

\[
  10^6 \times 32\ \text{octets} = 32 \times 10^6\ \text{octets} = \mathbf{32\ \text{Mo}}.
\]

%------------------------------------------------------------
\paragraph{Question 3 --- Capacité maximale avec 256 Mo.}

\[
  \frac{256 \times 10^6\ \text{octets}}{32\ \text{octets/nœud}} = \mathbf{8 \times 10^6}\ \textbf{positions au maximum.}
\]

\boitereponse{Dès que l'agent explore plus de 8 millions de positions, la RAM est saturée. \textbf{Une stratégie de remplacement est nécessaire} (ex. : LRU, ou remplacement aléatoire par profondeur).}

%------------------------------------------------------------
\paragraph{Question 4 --- Comparaison avec une table de hachage.}

Mémoire table de hachage : $2^{20} \times 12\ \text{octets} = 1\,048\,576 \times 12 \approx \mathbf{12\ \text{Mo}}$.

\begin{center}
\begin{tabular}{lcc}
\toprule
Critère                  & AVL              & Table de hachage \\
\midrule
Mémoire                  & 32 Mo (1M nœuds) & 12 Mo ($2^{20}$ entrées) \\
Coût d'accès (pire cas)  & $O(\log n)$      & $O(n)$ (collisions) \\
Coût d'accès (moyen)     & $O(\log n)$      & $O(1)$ \\
Taille                   & Dynamique        & Fixe \\
Gestion collisions       & Non              & Oui \\
Parcours ordonné         & Oui              & Non \\
\bottomrule
\end{tabular}
\end{center}

\boitereponse{%
La table de hachage est \textbf{plus rapide} ($O(1)$ moyen) et \textbf{moins consommatrice en mémoire} pour une taille fixe.
En revanche, sa taille est fixe (pas d'insertion illimitée) et les collisions dégradent les performances.
\textbf{En pratique pour une table de transposition de jeu, la table de hachage (type Zobrist) est préférée} car la vitesse d'accès prime sur le parcours ordonné des clés.%
}

% =====================================================================
\newpage
\section*{Exercice 3 --- Minimax et élagage alpha-bêta}
% =====================================================================

\subsection*{Partie A --- Minimax sur un arbre donné}

%------------------------------------------------------------
\paragraph{Question 1 --- Application de minimax.}

Structure de l'arbre : racine MAX (profondeur 0), niveau~1 MIN, niveau~2 MAX, feuilles au niveau~3.

Calcul de bas en haut :
\begin{itemize}
  \item \textbf{Niveau 2 (MAX)} :
    $C = \max(3, 12) = 12$, \quad
    $D = \max(8, 2) = 8$, \quad
    $E = \max(4, 6) = 6$, \quad
    $F = \max(14, 5) = 14$.
  \item \textbf{Niveau 1 (MIN)} :
    $A = \min(C, D) = \min(12, 8) = 8$, \quad
    $B = \min(E, F) = \min(6, 14) = 6$.
  \item \textbf{Racine (MAX)} :
    $\text{ROOT} = \max(A, B) = \max(8, 6) = \mathbf{8}$.
\end{itemize}

\begin{center}
\begin{tikzpicture}[
  level 1/.style={sibling distance=7cm},
  level 2/.style={sibling distance=3.5cm},
  level 3/.style={sibling distance=1.8cm},
]
  \node[maxnode]{8}
    child{node[minnode]{8}
      child{node[maxnode]{12}
        child{node[leafnode]{3}}
        child{node[leafnode]{12}}
      }
      child{node[maxnode]{8}
        child{node[leafnode]{8}}
        child{node[leafnode]{2}}
      }
    }
    child{node[minnode]{6}
      child{node[maxnode]{6}
        child{node[leafnode]{4}}
        child{node[leafnode]{6}}
      }
      child{node[maxnode]{14}
        child{node[leafnode]{14}}
        child{node[leafnode]{5}}
      }
    };
\end{tikzpicture}
\end{center}

\boitereponse{Valeur à la racine : \textbf{8}. Coup optimal pour MAX : aller vers le \textbf{fils gauche} (sous-arbre A, valeur 8).}

%------------------------------------------------------------
\paragraph{Question 2 --- Pseudo-code minimax.}

\begin{lstlisting}
int minimax(Node *node, bool maximizing) {
    if (is_leaf(node))           // cas de base : profondeur 0 ou etat terminal
        return node->value;

    if (maximizing) {
        int best = INT_MIN;
        for each child c of node:
            best = max(best, minimax(c, false));  // appel recursif MIN
        return best;
    } else {
        int best = INT_MAX;
        for each child c of node:
            best = min(best, minimax(c, true));   // appel recursif MAX
        return best;
    }
}
\end{lstlisting}

\textbf{Cas de base} : nœud feuille (profondeur $d = 0$ ou état terminal de jeu).\\
\textbf{Appels récursifs} : on alterne \texttt{maximizing = true/false} à chaque niveau.

%------------------------------------------------------------
\paragraph{Question 3 --- Complexité.}

À chaque nœud on explore $b$ enfants, sur $d$ niveaux :
\[
  \boxed{\text{Feuilles évaluées} = \Theta(b^d)}
\]

%------------------------------------------------------------
\paragraph{Question 4 --- Application Ataxx ($b = 40$, $d = 4$).}

\[
  40^4 = 2\,560\,000 \approx 2{,}56 \times 10^6\ \text{feuilles.}
\]

\boitereponse{C'est \textbf{faisable mais limité} : environ 2,5 millions d'évaluations par coup.
À $d = 5$ : $40^5 \approx 10^8$ évaluations $\to$ trop lent sans optimisation.
\textbf{Minimax pur seul n'est pas suffisant} pour des profondeurs utiles en pratique.}

% ============================
\subsection*{Partie B --- Élagage alpha-bêta}

%------------------------------------------------------------
\paragraph{Question 1 --- Déroulement pas à pas.}

Convention : $\alpha$ = meilleure valeur garantie pour MAX (borne inférieure courante),
$\beta$ = meilleure valeur garantie pour MIN (borne supérieure courante). Initialement $\alpha = -\infty$, $\beta = +\infty$.

\begin{enumerate}[label=\textbf{\arabic*.}]
  \item \textbf{ROOT}(MAX, $\alpha{=}{-}\infty$, $\beta{=}{+}\infty$) : descend vers A.
  \item \textbf{A}(MIN, $\alpha{=}{-}\infty$, $\beta{=}{+}\infty$) : descend vers C.
  \item \textbf{C}(MAX, $\alpha{=}{-}\infty$, $\beta{=}{+}\infty$) :
        évalue feuille 3 $\to$ $\alpha \leftarrow 3$ ;
        évalue feuille 12 $\to$ $\alpha \leftarrow 12$.
        Retourne \textbf{12}.
  \item \textbf{A} : $\beta \leftarrow \min(+\infty, 12) = 12$. Descend vers D.
  \item \textbf{D}(MAX, $\alpha{=}{-}\infty$, $\beta{=}12$) :
        évalue feuille 8 $\to$ $\alpha \leftarrow 8$ ($8 < 12$, pas de coupure) ;
        évalue feuille 2 $\to$ valeur = 8.
        Retourne \textbf{8}.
  \item \textbf{A} : $\beta \leftarrow \min(12, 8) = 8$. Retourne \textbf{8}.
  \item \textbf{ROOT} : $\alpha \leftarrow \max(-\infty, 8) = 8$. Descend vers B.
  \item \textbf{B}(MIN, $\alpha{=}8$, $\beta{=}{+}\infty$) : descend vers E.
  \item \textbf{E}(MAX, $\alpha{=}8$, $\beta{=}{+}\infty$) :
        évalue feuille 4 ($\alpha$ reste 8) ;
        évalue feuille 6 ($\alpha$ reste 8).
        Retourne \textbf{6}.
  \item \textbf{B} : $\beta \leftarrow \min(+\infty, 6) = 6$.
        Vérification : $\alpha \geq \beta$ ? $8 \geq 6$ \textbf{OUI}
        $\to$ \textcolor{red}{\textbf{coupure $\alpha$}} !
        F et ses feuilles (14, 5) ne sont \textbf{pas visités}.
  \item \textbf{ROOT} : retourne $\max(8, 6) = \mathbf{8}$.
\end{enumerate}

\begin{center}
\begin{tikzpicture}[
  level 1/.style={sibling distance=7cm},
  level 2/.style={sibling distance=3.5cm},
  level 3/.style={sibling distance=1.8cm},
]
  \node[maxnode]{8}
    child{node[minnode]{8}
      child{node[maxnode]{12}
        child{node[leafnode]{3}}
        child{node[leafnode]{12}}
      }
      child{node[maxnode]{8}
        child{node[leafnode]{8}}
        child{node[leafnode]{2}}
      }
    }
    child{node[minnode]{6}
      child{node[maxnode]{6}
        child{node[leafnode]{4}}
        child{node[leafnode]{6}}
      }
      child{node[prumednode]{\color{gris}---}
        child{node[prunedleaf]{\color{gris}14}}
        child{node[prunedleaf]{\color{gris}5}}
      }
    };
  % Annotation coupure
  \node[font=\small\bfseries, text=red] at (4.8, -3.5) {coupure $\alpha$ ici};
\end{tikzpicture}

\small Les nœuds en grisé/pointillés sont élaguées (non visités).
\end{center}

%------------------------------------------------------------
\paragraph{Question 2.}

\boitereponse{%
\textbf{Feuilles évaluées} : 3, 12, 8, 2, 4, 6 $\Rightarrow$ \textbf{6 feuilles}. \\
\textbf{Feuilles élaguées} : 14, 5 $\Rightarrow$ \textbf{2 feuilles}. \\
La valeur à la racine est \textbf{inchangée : 8}.%
}

%------------------------------------------------------------
\paragraph{Question 4 --- Rôle de $\alpha$ et $\beta$.}

\boitereponse{%
$\alpha$ est la meilleure valeur que MAX est \emph{déjà assuré} d'obtenir sur un autre chemin (borne inférieure).
$\beta$ est la meilleure valeur que MIN est \emph{déjà assuré} de limiter (borne supérieure).

Si $\alpha \geq \beta$ : MAX a déjà un chemin valant $\geq \alpha$, mais dans le sous-arbre courant MIN va imposer une valeur $\leq \beta < \alpha$. MAX ne choisira jamais ce sous-arbre : \textbf{on peut l'élaguer sans perte d'information}.%
}

%------------------------------------------------------------
\paragraph{Question 5 --- Gain en profondeur effective.}

\begin{itemize}
  \item \textbf{Minimax} : $40^4 = 2\,560\,000$ feuilles.
  \item \textbf{Alpha-bêta meilleur cas} : $b^{d/2} = 40^2 = \mathbf{1\,600}$ feuilles.
  \item \textbf{Gain} : $\dfrac{40^4}{40^2} = 40^2 = 1600$ fois moins de feuilles évaluées.
\end{itemize}

\boitereponse{%
Pour le même budget de calcul, alpha-bêta (coups bien ordonnés) peut explorer jusqu'à \textbf{profondeur $d = 8$}
là où minimax ne peut aller qu'à $d = 4$, car $40^{8/2} = 40^4$.
\textbf{Alpha-bêta double la profondeur d'exploration effective.}%
}

% ============================
\subsection*{Partie C --- Ordonnancement et table de transposition}

%------------------------------------------------------------
\paragraph{Question 1 --- Pire cas alpha-bêta.}

\boitereponse{%
Dans le \textbf{pire ordre} (coups explorés du moins bon au meilleur), aucune coupure n'est déclenchée.
Alpha-bêta évalue \textbf{$b^d$ feuilles}, identique au minimax pur.%
}

%------------------------------------------------------------
\paragraph{Question 2 --- Critères d'ordonnancement des coups.}

\begin{enumerate}
  \item \textbf{Nombre de captures} : trier les coups par ordre décroissant du nombre de pièces adverses retournées. Les coups offensifs sont statistiquement meilleurs et génèrent des coupures plus tôt.
  \item \textbf{Score de la table de transposition} : si la position résultant d'un coup a déjà été évaluée (hash Zobrist présent dans l'AVL), utiliser ce score pour ordonner les coups du plus prometteur au moins bon.
\end{enumerate}

%------------------------------------------------------------
\paragraph{Question 3 --- Rôle de la table de transposition (AVL de l'exercice 2).}

\textbf{(a) Éviter les réévaluations :}
Avant chaque appel récursif, on calcule le hash Zobrist de l'état courant et on le cherche dans l'AVL.
Si le hash est présent avec une profondeur d'analyse $\geq$ la profondeur demandée, on retourne directement le score stocké sans descendre dans l'arbre de jeu.

\textbf{(b) Améliorer l'ordonnancement :}
Pour chaque coup légal, on calcule le hash de la position résultante et on consulte l'AVL.
Si un score est disponible, il sert de clé de tri pour explorer les meilleurs coups en premier, maximisant ainsi les coupures alpha-bêta.

%------------------------------------------------------------
\paragraph{Question 4 --- Tableau de synthèse.}

\begin{center}
\renewcommand{\arraystretch}{1.4}
\begin{tabular}{lccc}
\toprule
\textbf{Algorithme} & \textbf{Feuilles évaluées} & \textbf{Profondeur ($b{=}40$)} & \textbf{Table transp.} \\
\midrule
Minimax pur                   & $b^d$              & $d \leq 3$    & Non \\
Alpha-bêta (pire cas)         & $b^d$              & $d \leq 3$    & Non \\
Alpha-bêta (meilleur cas)     & $b^{d/2}$          & $d \leq 6$    & Non \\
Alpha-bêta + ordonn. + TT     & $\approx b^{d/2}$  & $d \leq 6$--$8$ & Oui \\
\bottomrule
\end{tabular}
\end{center}

\textit{Note} : seuil pratique estimé à $\approx 10^6$ feuilles par coup pour un temps de réponse acceptable.\\
Minimax : $40^3 = 64\,000$ (\checkmark), $40^4 = 2{,}56M$ (limite).
Alpha-bêta meilleur cas : $40^3 \Rightarrow d = 6$ ; $40^4 \Rightarrow d = 8$.

\end{document}
